\chapter{Lecture 03}

\hfill {\footnotesize \textit{Date: 13/10/2026}}

\section{Review on linear algebra---continued}

\subsection{Orthogonality}

\begin{definition}[Orthogonality]
	Let \( V \) be an inner product space. Two vectors \( \mathbf{x}, \mathbf{y} \in V \) are \emph{orthogonal} if
	\begin{equation}
		\langle \mathbf{x}, \mathbf{y} \rangle = 0.
		\label{eq:orthogonality-1}
	\end{equation}
\end{definition}



\begin{example}
	Let \( \mathbf{x} = \begin{bmatrix} 1 \\ 2 \\ 1 \end{bmatrix}, \mathbf{y} = \begin{bmatrix} 1 \\ -1 \\ 1 \end{bmatrix} \in \mathbb{R}^3 \). Then

	\begin{equation}
		\langle \mathbf{x}, \mathbf{y} \rangle = 1 \times 1 + 2 \times (-1) + 1 \times 1 = 0.
		\label{eq:orthogonality-2}
	\end{equation}

	Therefore, \( \mathbf{x}, \mathbf{y} \) are orthogonal/perpendicular.

	In picture,
\end{example}
\todo{Add 3d picture of these orthogonal vectors.}

But since inner product is defined for any \( \mathbb{R}^n \), we can talk about orthogonality there as well.

\begin{example}
	Let \( \mathbf{x} = \begin{bmatrix} 1 \\ 2 \\ 1 \\ 0 \\ -7 \end{bmatrix}, \mathbf{y} = \begin{bmatrix} 1 \\ -1 \\ 1 \\ 5 \\ 0 \end{bmatrix} \in \mathbb{R}^3 \). Then

	\begin{equation}
		\langle \mathbf{x}, \mathbf{y} \rangle = 1 \times 1 + 2 \times (-1) + 1 \times 1 + 0 \times 5 + (-7) \times 0 = 0.
		\label{eq:orthogonality-2.1}
	\end{equation}

	Therefore, \( \mathbf{x}, \mathbf{y} \) are orthogonal.

	Notice, \( \mathbf{x}, \mathbf{y} \) are vectors in \( 5 \)-dimensional space  \( \mathbb{R}^5 \), which we can not visualise, yet we can talk about whether they are orthogonal/perpendicular to each other using the language of (linear) algebra.
\end{example}

If the basis vectors $v_1,v_2,\ldots,v_n$ of the vector
space $\mathbb{R}^n$ are pairwise orthogonal, i.e.,
\begin{equation}
	\langle v_i,v_j\rangle=0,
	\qquad i\neq j,
	\label{eq:orthogonality-3}
\end{equation}
then the coefficients of any vector $v\in\mathbb{R}^n$
with respect to this basis can be determined easily.
Write
\begin{equation}
	v=\sum_{i=1}^{n}c_i v_i,
\end{equation}
and taking the inner product with $v_j$ on both sides,
we obtain
\begin{align}
	\langle v,v_j\rangle
	 & =\left\langle\sum_{i=1}^{n}c_i v_i,v_j\right\rangle \\
	 & =\sum_{i=1}^{n}c_i\langle v_i,v_j\rangle            \\
	 & =c_j\langle v_j,v_j\rangle,
\end{align}
since $\langle v_i,v_j\rangle=0$ whenever $i\neq j$.
Therefore,
\begin{equation}
	c_j=\frac{\langle v,v_j\rangle}
	{\langle v_j,v_j\rangle}
	=\frac{\langle v,v_j\rangle}{\|v_j\|^2}
	\label{eq:orthogonality-4}
\end{equation}
Thus, orthogonality allows us to determine each
coefficient independently, without solving a system
of linear equations. In particular, if the basis
is orthonormal, then
\begin{equation}
	v_j=\langle v,v_j\rangle
	\label{eq:orthogonality-5}
\end{equation}

\begin{example}

\end{example}
\todo{Do an example in \( \mathbb{R}^3 \) is possible, otherwise,}

Let \( V \) be a normed space.

\begin{definition}[Normalised vectors]
	A vector \( v \in V \) is called \emph{normalised} if
	\begin{equation}
		\|v\| = 1.
		\label{eq:normalised-vectors-1}
	\end{equation}
\end{definition}

Suppose \( v \in V \) is not normalised. But
\begin{equation}
	\frac{v}{\|v\|}
	\label{eq:orthogonality-6}
\end{equation}

is normalised. Because

\begin{equation}
	\left\|\frac{v}{\|v\|}\right\| =\footnote{Using absolute homogeneity property of norm} \frac{1}{\|v\|} \|v\| = 1.
	\label{eq:orthogonality-7}
\end{equation}

The process making a vector normalised in this way, is called \emph{normalisation}.

\begin{definition}[Orthogonal basis]
	A set of basis \( \{v_1, v_2, \ldots, v_n\} \) of a vector space \( V \) is called \emph{orthogonal basis} if
	\begin{equation}
		\langle v_i, v_j \rangle = \begin{cases}
			c \in \mathbb{R} & \text{if } i = j     \\
			0                & \text{if } i \neq j.
		\end{cases}
		\label{eq:orthogonal-basis-1}
	\end{equation}
\end{definition}

\begin{definition}[Orthonormal basis]
	A set of basis \( \{v_1, v_2, \ldots, v_n\} \) of a vector space \( V \) is called \emph{orthonormal basis} if
	\begin{equation}
		\langle v_i, v_j \rangle = \begin{cases}
			1 & \text{if } i = j     \\
			0 & \text{if } i \neq j.
		\end{cases}
		\label{eq:orthogonal-basis-2}
	\end{equation}
\end{definition}

\textbf{Notation (Kronecker/Dirac delta).}
\begin{equation}
	\delta_{ij} = \begin{cases}
		1, i = j \\
		0, i \neq j.
	\end{cases}
	\label{eq:orthogonality-8}
\end{equation}

\todo{Add examples related to Dirac delta}

Now we prove \autoref{thm:eigsofsymmatrix}.

\begin{proof}[Proof of \autoref{thm:eigsofsymmatrix}]
	\todo{Complete the proof}
\end{proof}

In addition,

\begin{theorem}
	Eigenvectors with different eigenvalues of a symmetric matrix are orthogonal.
\end{theorem}
\begin{proof}
	\todo{Complete the proof.}
\end{proof}

\todo{Construction of symmetric matrices}

\begin{definition}[Anti/skew symmetric matrix]
	A matrix \( A \) is called  \emph{anti/skew symmetric} if
	\begin{equation}
		A = -A^T.
		\label{eq:anti-skew-symmetric-matrix-1}
	\end{equation}
\end{definition}

\todo{Construction of anti-symmetric matrices}

Most matrices are neither symmetric nor anti symmetric. But given a matrix \( M \), it can be written as sum of one symmetric and one anti symmetric matrix as follows:

\todo{Complete}

\section{Review on multivariable calculus}

In this course, we will optimise functions whose inputs and outputs may consist of several real numbers. A function with several input variables is called a
\emph{multivariable function}. The branch of mathematics concerned with such functions is known as \emph{multivariable calculus}.

\subsection{Terminology Related to Functions}

\begin{definition}
	Let $A$ and $B$ be nonempty sets. A function $f:A\to B$ is a rule that assigns to every element $x\in A$ exactly one element $f(x)\in B$. Here,

	\begin{itemize}
		\item \( A \) is called the \emph{domain} of $f$.

		\item \( B \) is called the \emph{codomain} of $f$.

		\item The \emph{range}, or \emph{image}, of $f$ is the subset of $B$ consisting of all values actually attained by $f$, namely $f(A)=\{f(x):x\in A\}$.

		\item The \emph{graph} of $f$ is the set $\operatorname{Graph}(f) =\{(x,f(x)):x\in A\}$.
	\end{itemize}
\end{definition}

\begin{example}
	Consider the function $f:\mathbb{R}\to\mathbb{R}$ defined by $f(x)=x^2$.

	Its domain and codomain are both $\mathbb{R}$, but range is \( [0, \infty) \).

	The graph of $f$ consists of all points of the form $(x,x^2)$, where $x\in\mathbb{R}$. A portion of this graph is shown below.
\end{example}

\todo{add grpah of \( f(x) = x^2 \)}

Graph of a function \( f: \mathbb{R} \to \mathbb{R} \) can be represented in a plane. But for functions of several variables, the situation
is more complicated. For example, the graph of a function $f:\mathbb{R}^2\to\mathbb{R}$ lies in $\mathbb{R}^3$ and can often be visualized as a surface. However, the graph of a function $f:\mathbb{R}^2\to\mathbb{R}^2$ lies in $\mathbb{R}^4$, which cannot be represented directly in ordinary three-dimensional space.

\todo{Add graph of \( x^2 + y \)}

\subsection{Directional Derivatives}

\subsubsection*{The Derivative of a Single-Variable Function}

We first recall the notion of differentiation from elementary calculus.

\begin{definition}
	Let $I\subseteq\mathbb{R}$ be an open interval, let $f:I\to\mathbb{R}$, and	let $a\in I$. We say that $f$ is \emph{differentiable at $a$} if the following limit exists as a finite real number:
	\begin{equation}
		f'(a)
		=
		\lim_{h\to 0}
		\frac{f(a+h)-f(a)}{h}.
		\label{eq:single-variable-derivative}
	\end{equation}
	The value $f'(a)$ is called the \emph{derivative} of $f$ at $a$.
\end{definition}

Geometrically, the quotient
$\frac{f(a+h)-f(a)}{h}$ represents
the slope of the secant line joining
$(a,f(a))$ and $(a+h,f(a+h))$.
As $h$ approaches zero, the secant
slope approaches the slope of the
tangent line, provided the limit exists.

Suppose now that
$f:U\to\mathbb{R}$, where
$U$ is an open subset of
$\mathbb{R}^n$ and $n\geq 2$.

Given $a\in U$, a small change
in the input is represented by a
vector $h\in\mathbb{R}^n$, rather
than by a real number. One might
attempt to imitate
\eqref{eq:single-variable-derivative}
by considering
$\frac{f(a+h)-f(a)}{h}$.

However, this expression is not
defined in the usual arithmetic
of Euclidean spaces: its numerator
is a real number, whereas its
denominator is a vector. There
is no canonical operation of
division by a nonzero vector
analogous to division by a
nonzero real number.

There is also a geometric difficulty.
In $\mathbb{R}$, a point can be
approached only from the left or
from the right. In $\mathbb{R}^n$,
with $n\geq 2$, a point may be
approached along infinitely many
different directions. The rate
of change of a function may
depend on the particular direction
in which its input is varied.

This leads us to the following definition:

\begin{definition}
	Let $U\subseteq\mathbb{R}^n$
	be open, let
	$f:U\to\mathbb{R}^m$,
	and let $a\in U$.
	For a nonzero vector
	$v\in\mathbb{R}^n$, the
	\emph{directional derivative}
	of $f$ at $a$ in the direction
	$v$ is defined by
	\begin{equation}
		D_vf(a)
		=
		\lim_{t\to 0}
		\frac{f(a+tv)-f(a)}{t},
		\label{eq:directional-derivative}
	\end{equation}
	provided that this limit exists
	in $\mathbb{R}^m$.
\end{definition}

The directional derivative
$D_vf(a)$ measures the
instantaneous rate of change
of $f$ at $a$ as the input
moves along the line
$a+tv$. When $v$ is a unit
vector, it measures the
rate of change per unit
distance in that direction.

\begin{remark}
	The ordinary derivative
	is a special case of the
	directional derivative.
	Indeed, when $n=m=1$
	and $v=1$, the definition
	reduces to
	$D_1f(a)=f'(a)$.
\end{remark}
